\section{Discussion}

Obwohl eine prokaitve KI vereinzelt zu neuen Perspektiven innerhlab des Meeting führt und tatsächlic in
eingen fälle den Teilnehmden geholfen hat die Redeantiele fairer zu verteielte und auf Themen aufmerksam
gemacht hat die sonst untergegagen wären, überweigt in diesem System entwurf der Negative impakt. Nahe zu
alle Teilnehmden empfanden die KI als störend oder sahen, dass die KI sich zu oft beteiligt.

diese nachteiel überweigen den geringen positiven effekten, welche das system auf das meeting hatte.

Die shclussfolgernung, dass ein proaktiver KI-Faciliatoar generell nicht sinnvoll ist, wäre allerdings zuverführt.

Druch anpassen der Therholds könnte die interaktions rate der KI veringert werden und auf die meeting
situation angepast werden.

Der reltiv navier Ansatz die KI selbst ihre gendaen mit einer genreischen Bewertunglogik zu bewerten, fürhte
zu generischen Antworten mit immer ähnlichen pattern. (Erst bedanekn, dann paraphrasieren) dies ist zwar aus
den schilderunge der experten interviews modleiert, wruden dann dennoch in der erhöhte frequenz als nachteil empfunden.

Ein PRoaktiver ansazt benötigt deutlich mehr genauigkeit. im zeifel wäre es vermutlich besser gewsen, wenn
die KI nur punktuell eingegirffen hätte. anstatt permant neue Gedanken zu formen udn zu evaluieren.
Ein die geformten gednekn hatten ohnehin eine hohe ähnlichkeit, weshalb der merhwert im faciliatoatr zenario
infrage gestellt werden kann.

Aus den beobachtungen diese Prototypen udn der user studies leitet der Autor folgende empfheungen ab:

DEr versuchsaubgau war gut geeiengt um ein solches system unter live bedinugnen zu testen.
allerdings hätten davor erstmal einzelen funktionen isoliert geteste werden sollen.
Ein navierer ansatz und vermutlich zielf+hrender wäre z.B. nur den context workflow laufen zulassen. und die
inner thougs komplett wegzulassen.
dafür hätte es noch die notwendigkeit gegeben, den ki agenent zu unterbrechen.

% ---------------------------------------------------------------------------------------------------------------------
\subsection{Lessons Learned: Implications for similar System Designs}

Im druchschnitt hat die KI alle x minuten interveniwert. man könnte dies ggf. darduch reduzieren, dass der
interrupt und der imThreshold heruntergesetzt werden. Allerdings würde dies dazuführen, dass die meisten
gedanekn unausgesprochen bleiben. Dies wirk erstmal verschwenderisch. es müsste eine gleichgewicht gefunden
werden dass wirklich relevante gedanekn in dem aktuellen moment angepsrochen werden. andere wiederum

Besser kontext, wenn eien genaua agend vorliegt und facilitatorin anweiseungen was in diesm agenda puntk
besonders relevant ist. das könnten die antworten verbesseren. Hinzukommt auch eine grnularere ein teilung in
phasen. in diesem beipsiel hätte z.B. der Facilitator am anfang ein vorgehen vorschalgen können. "Wir sammeln
Reihum die vor und nachteile der kandidaten und entscheiden dann welcher kandidat am besten passt."

Zukünfige arbeiten sollten nicht versuchen einen generischen Facilaitor zu entwciklen, welche passend
reagiert, sondern sich bewährte strukturen nutzen diesn ablauf festlegen, dann die eingschaften und wichtigen
Punkte für die Phasen dieser strukur ableiten und klar formulieren. ein KI-Facilitator würde sich dann auf.
somit kann dann besser evaluiert werden, ob der facilitator sich an die merkamle dieser phasen gehalten hat.
die reduziert die aufgaben des fecilaitoas auf wenige haupt punkt und würde vermutlich einen größern impakt erzeugen.

Es wäre dann möglihc die heurisiken desinner thoughs models an diese granualren phasen anzupassen.

[TODO: Abbildung mit phasen, wie ein weiteres System design aussehen könnte.]

Im Zuge dessen wäre es dann entscheident das logging der LLM-Traces zu verbessern, um das system besser
evaluieren zu können.

entscheidender währe dann die übergänge zwischnen den phasen.

Es gibt neben effekte die das evaluaieren des systems erschweren: die trun allocation könnten
fälschlicherweise den faciliatoatr zu weisen, dann reicht der imThreshoad um einen gedanken zu äußern. da
sich die gedanekn aber warhshcienlcih eh, um ein relevantes theam drehen führt das nicht zwangsläufig zu
eienr falschen intervention. Sei kann trotzdem passend sein. (ggf. Beispiele ergänzen)

gedanken die geäußert wurden , waren immer sehr ähnlich. und die auswertung udn das monitoring von sovielen
traces ist sehr aufwendig. vereinfachen und gesagtes buffern udn dann in einem
batch verarbieten.

Obwohl das problem des overtalking im system design erkennt wurde und druch dei ageing funktion veruscht
wrude das äußer des selben gendekn zu ferindern, war dies nicht aussreichend. Vermutlich hätte eine cooldown
auf die sprache generell die qualität verbessert. Zwar hätte die ki hier relevante informatione verpassen
könenn, allerdings wäre dies ein tradeoff welche hinnehmbar wäre, wenn dadruch die akzeptanz steigen würde

Live interaktion utnerleigen einen höhren anspruch als chat interaktionen. das inner thogs framewokr wurde
als chat framework konzipiert. in chats ist es ünlich das eine nachricht auchmal nicht ganz genau auf den
aktuelln kontext passt. (user sind es geäht, dass das tippen dauert und es sit eien asyncrone interaktione).
bei live interaktione ist die ander. teilenhmden fallen latzen auf. wenn nicht direkt informationen

Tracebility ist aufwändig und schwer zu interpretieren. eine platform wie langfuse hat während der
Entwikclung enorm geholfen ein intuitives verständniss für das system zu entwikclen.
Jedoch sollte sich gut gedanken gemacht werden wie die
traces in einer solch platform gelogget werden, um sie später zu analysieren.

bereits nach der ersten interation des Systems wruden deutlich schwachstellen ersichtlich. der besser ansatz
wäre gewesen das system direkt zu verbessern und dann die nächste user gruppe drauf zu setzen. für diese
prototpy evaluierung wären kleiner phasen ggf besser gewesen.